\documentclass[12pt,a4paper]{ctexart}
\usepackage[left=2.5cm,right=2.5cm,top=2.3cm,bottom=2.3cm]{geometry}
\usepackage{xcolor}
\usepackage{tcolorbox}
\usepackage{titlesec}
\usepackage{fancyhdr}
\usepackage{enumitem}
\usepackage{tabularx}
\usepackage{amsmath,amssymb}
\usepackage{hyperref}
\tcbuselibrary{skins,breakable}

\definecolor{MainColor}{HTML}{2A6F73}
\definecolor{SecondColor}{HTML}{3CAEA3}
\definecolor{LightMain}{HTML}{EAF7F6}
\definecolor{TitleDark}{HTML}{1F3A40}
\definecolor{TextGray}{HTML}{555555}
\definecolor{BorderGray}{HTML}{CBD5D6}
\definecolor{WarnColor}{HTML}{F2B35D}
\definecolor{ErrorColor}{HTML}{C94C4C}
\definecolor{DefineColor}{HTML}{6C63B7}

\linespread{1.25}
\setlength{\parindent}{2em}
\setlength{\parskip}{0.25em}
\hypersetup{colorlinks=true,linkcolor=MainColor,urlcolor=MainColor}
\pagestyle{fancy}
\fancyhf{}
\fancyfoot[L]{\textcolor{MainColor}{机器学习笔记}}
\fancyfoot[R]{\textcolor{MainColor}{\thepage}}
\renewcommand{\headrulewidth}{0pt}
\renewcommand{\footrulewidth}{0pt}

\newcommand{\topbar}[2]{\noindent{\small\textcolor{MainColor}{#1}}\hfill{\small\bfseries\textcolor{TitleDark}{#2}}\par\vspace{0.35em}{\color{BorderGray}\hrule}\vspace{1.2em}}
\newcommand{\notetitlecard}[6]{
\begin{tcolorbox}[enhanced,colback=white,colframe=BorderGray,boxrule=0.6pt,arc=2mm,left=12pt,right=12pt,top=10pt,bottom=10pt,borderline west={4pt}{0pt}{SecondColor}]
\begin{tabularx}{\linewidth}{X r}
{\small\textcolor{TextGray}{#1}} & {\small\textcolor{TextGray}{#2}} \\[0.9em]
\multicolumn{2}{l}{{\Huge\bfseries\textcolor{MainColor}{#3}}} \\[0.45em]
\multicolumn{2}{l}{{\large\textcolor{SecondColor!90!black}{#4}}} \\[0.45em]
\multicolumn{2}{l}{{\small\textcolor{TextGray}{课程：#5 \qquad 作者：#6}}}
\end{tabularx}\end{tcolorbox}\vspace{1em}}
\newtcolorbox{summarybox}[1][本节摘要]{enhanced,breakable,colback=LightMain,colframe=SecondColor,colbacktitle=SecondColor,coltitle=white,title={#1},fonttitle=\bfseries,boxrule=0.6pt,arc=2mm,left=12pt,right=12pt,top=10pt,bottom=10pt}
\newtcolorbox{definitionbox}[1][定义]{enhanced,breakable,colback=DefineColor!6,colframe=DefineColor,colbacktitle=DefineColor,coltitle=white,title={#1},fonttitle=\bfseries,boxrule=0.6pt,arc=2mm,left=12pt,right=12pt,top=10pt,bottom=10pt}
\newtcolorbox{keybox}[1][重点]{enhanced,breakable,colback=MainColor!5,colframe=MainColor,colbacktitle=MainColor,coltitle=white,title={#1},fonttitle=\bfseries,boxrule=0.6pt,arc=2mm,left=12pt,right=12pt,top=10pt,bottom=10pt}
\newtcolorbox{examplebox}[1][例题]{enhanced,breakable,colback=white,colframe=SecondColor,colbacktitle=SecondColor!85!black,coltitle=white,title={#1},fonttitle=\bfseries,boxrule=0.6pt,arc=2mm,left=12pt,right=12pt,top=10pt,bottom=10pt}
\newtcolorbox{mistakebox}[1][易错点]{enhanced,breakable,colback=ErrorColor!5,colframe=ErrorColor,colbacktitle=ErrorColor,coltitle=white,title={#1},fonttitle=\bfseries,boxrule=0.6pt,arc=2mm,left=12pt,right=12pt,top=10pt,bottom=10pt}
\newtcolorbox{tipbox}[1][提示]{enhanced,breakable,colback=WarnColor!10,colframe=WarnColor,colbacktitle=WarnColor,coltitle=white,title={#1},fonttitle=\bfseries,boxrule=0.6pt,arc=2mm,left=12pt,right=12pt,top=10pt,bottom=10pt}
\titleformat{\section}{\Large\bfseries\color{MainColor}}{\thesection}{0.8em}{}
\titleformat{\subsection}{\large\bfseries\color{TitleDark}}{\thesubsection}{0.8em}{}
\titleformat{\subsubsection}{\normalsize\bfseries\color{SecondColor!80!black}}{\thesubsubsection}{0.8em}{}
\setlist[itemize]{leftmargin=2.2em,itemsep=0.3em,topsep=0.4em}
\setlist[enumerate]{leftmargin=2.2em,itemsep=0.3em,topsep=0.4em}
\newcommand{\important}[1]{\textbf{\textcolor{MainColor}{#1}}}
\newcommand{\warning}[1]{\textbf{\textcolor{ErrorColor}{#1}}}
\newcommand{\concept}[1]{\textbf{\textcolor{DefineColor}{#1}}}

\begin{document}
\topbar{吴恩达机器学习}{学习笔记}
\notetitlecard{Study Notes Series}{2026 Spring}{吴恩达机器学习笔记(5)}{}{机器学习}{C++与草稿纸}

\begin{summarybox}
本周讨论如何训练第四周定义的神经网络：用代价函数衡量输出误差，用反向传播高效计算每个权重的梯度，再用梯度下降或高级优化器更新参数。还介绍展开参数、梯度检验、随机初始化和完整训练流程。
\begin{itemize}[leftmargin=2em,itemsep=2pt,topsep=4pt]
  \item 多类神经网络使用交叉熵代价，并对非偏置权重进行正则化；
  \item 反向传播从输出层向输入层传播误差，得到每个 $\Theta^{(l)}$ 的梯度；
  \item 梯度检验用数值导数核对反向传播实现，随机初始化打破对称性；
  \item 自主驾驶等系统通常由感知、预测和控制等多个学习模块组成。
\end{itemize}
\end{summarybox}

\section{神经网络的代价函数}
\begin{definitionbox}[多类分类代价]
设网络有 $K$ 个输出单元，样本标签用 one-hot 向量 $\mathbf y^{(i)}$ 表示，输出为 $\mathbf a^{(L,i)}$。多类交叉熵代价可写为
\[
J(\Theta)=-\frac1m\sum_{i=1}^{m}\sum_{k=1}^{K}
\left[y_k^{(i)}\log a_k^{(L,i)}+(1-y_k^{(i)})\log(1-a_k^{(L,i)})\right].
\]
对所有训练样本和输出单元求平均，模型训练就是寻找使 $J(\Theta)$ 最小的权重矩阵集合。
\end{definitionbox}
\begin{keybox}[正则化神经网络]
为抑制过拟合，在代价函数中加入除偏置列以外的权重平方和：
\[
J_{\mathrm{reg}}(\Theta)=J(\Theta)+\frac{\lambda}{2m}
\sum_{l=1}^{L-1}\sum_{i}\sum_{j=1}^{s_{l+1}}\left(\Theta_{ji}^{(l)}\right)^2.
\]
第一列通常对应偏置权重，不参与正则化。$\lambda$ 越大，模型越倾向于使用较小的权重。
\end{keybox}

\section{反向传播算法}
\begin{definitionbox}[误差项]
对第 $i$ 个样本，记第 $l$ 层的误差为 $\delta^{(l)}$。输出层误差由预测与标签之差给出：
\[
\delta^{(L)}=\mathbf a^{(L)}-\mathbf y.
\]
对隐藏层，误差从下一层反向传播：
\[
\delta^{(l)}=\left((\Theta^{(l)})^T\delta^{(l+1)}\right)\odot
g'\!\left(\mathbf z^{(l)}\right),
\qquad l=L-1,\ldots,2,
\]
其中 $\odot$ 表示逐元素乘法。偏置单元对应的误差不参与下一层传播。
\end{definitionbox}
\begin{keybox}[反向传播流程]
对每个训练样本执行：
\begin{enumerate}
  \item 前向传播，保存每层 $\mathbf z^{(l)}$ 和 $\mathbf a^{(l)}$；
  \item 计算输出层误差 $\delta^{(L)}$；
  \item 从 $L-1$ 层向第 $2$ 层反向计算误差；
  \item 累加梯度：
  \[
  \Delta^{(l)}\leftarrow\Delta^{(l)}+\delta^{(l+1)}(\mathbf a^{(l)})^T;
  \]
  \item 取平均并加上正则化项，得到 $\partial J/\partial\Theta^{(l)}$。
\end{enumerate}
反向传播的意义是利用链式法则共享中间结果，避免对每个参数重复计算完整的数值导数。
\end{keybox}
\begin{examplebox}[Sigmoid 导数]
若 $g(z)=1/(1+e^{-z})$，则
\[
g'(z)=g(z)(1-g(z)).
\]
因此保存 $a^{(l)}=g(z^{(l)})$ 后，可以直接用 $a^{(l)}\odot(1-a^{(l)})$ 计算 Sigmoid 的导数。
\end{examplebox}

\section{实现与调试}
\subsection{展开参数}
\begin{definitionbox}[把矩阵参数变成向量]
数值优化器通常要求参数是一个长向量。实现时把所有 $\Theta^{(l)}$ 按固定顺序拼接成 $\theta$，计算代价前再恢复为矩阵。展开和恢复必须使用完全一致的顺序，否则代价和梯度会对应到错误的参数。
\end{definitionbox}
\subsection{梯度检验}
\begin{tipbox}[数值梯度检查]
对第 $j$ 个参数，用中心差分近似导数：
\[
\frac{\partial J}{\partial\theta_j}\approx
\frac{J(\theta+\varepsilon\mathbf e_j)-J(\theta-\varepsilon\mathbf e_j)}{2\varepsilon}.
\]
将该数值结果与反向传播的解析梯度逐项比较。梯度检验只用于小规模调试，确认正确后必须关闭，否则计算成本会非常高。
\end{tipbox}
\begin{mistakebox}[反向传播的常见错误]
\begin{itemize}
  \item 把偏置列也加入正则化；
  \item 忘记误差传播中的转置或逐元素乘法；
  \item 梯度检验时仍使用随机变化的参数、数据或 dropout；
  \item 只检查代价不检查梯度，导致优化器无法正常下降；
  \item 用过大的学习率，使代价函数发散。
\end{itemize}
\end{mistakebox}

\subsection{随机初始化与完整流程}
\begin{keybox}[随机初始化的作用]
如果所有权重都初始化为相同值，同一层的神经元会得到完全相同的梯度，训练后仍然保持相同功能，这称为对称性问题。通常从小范围均匀分布中随机初始化权重，例如
\[
\Theta_{ij}^{(l)}\sim U(-\epsilon,\epsilon).
\]
初始化范围应与上一层和下一层的规模相称。
\end{keybox}
\begin{tipbox}[训练神经网络的工程流程]
确定层数和每层单元数，随机初始化参数；实现前向传播和代价函数；实现反向传播得到梯度；用梯度检验验证实现；再调用优化器最小化正则化代价；最后在验证集和测试集上检查泛化能力。
\end{tipbox}

\section{应用：自主驾驶}
\begin{examplebox}[从感知到控制]
自主驾驶系统可以把摄像头图像、雷达和其他传感器作为输入，通过神经网络识别道路、车辆、行人和车道线，再由规划与控制模块决定转向、加速和制动。实际工程中常将系统拆为多个可单独评估的模块，并用大量覆盖边界情况的数据验证每个模块。
\end{examplebox}
\section{本周知识脉络}
\begin{keybox}[从前向传播到反向传播]
前向传播计算预测，代价函数衡量预测与标签的差异；反向传播利用链式法则计算所有权重的梯度；正则化控制模型复杂度；梯度检验、随机初始化和分模块评估保证实现可以可靠训练。下一步应结合偏差、方差和验证集选择网络结构与正则化强度。
\end{keybox}
\end{document}
